\begin{abstract}
在 2026 年大模型神仙打架的语境下，部署全员由 GPT-5.3、Gemini 3.1 Pro 或 Claude Sonnet 4.6 组成的 SOTA 多智能体系统，不仅在 API 开销上令人望而却步，也违背了工程学上的资源最优配置原则。现实中，广大研究者和开发者手里往往只有大批价格极其低廉、但也常常伴随着“智商堪忧”、工具调用格式不稳定、且极爱引发幻觉级联的次级模型（如 MiniMax 系列或 Flash-Lite 架构）。如何让这群“不靠谱”的廉价模型参与到高度严谨的科研推理中？本文提出的 \textbf{Grox\_chat} (Gemini Research Orchestration with minimaX) 框架，给出了一套极具“黑色幽默”与工程硬核感的解决方案。

我们没有奢望去提高廉价模型的代码执行成功率，也没有赋予它们危险的 OpenClaw 沙盒权限，而是把它们关进了一个由 LangGraph 状态机驱动、规则极其严苛的物理“监狱”。Grox\_chat 的核心创新在于“去中心化的物理牵制”：(1) **自然语言意图 RAG**，彻底褫夺了廉价模型的 JSON 工具调用权，只允许它们输出纯中文检索需求，由底层系统代为执行向量化和 Cross-Encoder 精排；(2) 外部挂载的 **SQLite 持久化黑板**，完全剥夺了模型自己维护长上下文的资格，彻底根治了多轮对话中的遗忘症；(3) 创造性地引入了游离于常规辩论之外的**“非理性气氛组”——狗 (Dot)、猫 (Cat) 与裁决者 (Tron)**。这组纯粹为了找茬、护主与维护系统安全法则而存在的节点，通过在每一轮末尾强行插入惩罚或奖励回合（DAPA路由机制），竟逼迫满嘴跑火车的专家模型在互相撕咬中实现了惊人的自我纠偏。实验证明，只要“系统物理牵制”足够强硬，即便是最劣质的基座模型，也能在 Grox\_chat 的 Chat-Only 基线下，被压榨出足以匹敌高端旗舰模型的逻辑推理质量。
\end{abstract}